\section{Limitations, Future Work, and Broader Impact}
The model optimizes a declared workload, finite implementation library, and
analytical hardware contract. Full-library pricing and global master optimality
remain open. Continuous relaxation bounds, feasible schedules, numerical
inference equivalence, and physical timing are distinct claims. The full-input
DeepSeek result uses captured routes and expanded attention under its recorded
precision policy; its workload is reproducible, not statistically representative
of every deployment. The large speedup depends on the reference and custom
memory-interface assumptions. Same-chip controls isolate software effects but
are not globally reoptimized ablations.

The current execution credits compatible batching, KV retention/placement,
and SwiGLU storage fusion. Larger native fusion groups, unrestricted overlap,
quality-changing menus, and general rematerialization gains need additional
experiments. RTL and synthesis studies cover selected schedules and datapaths,
not complete-chip calibration. Extending the framework to multiple chips,
robust workload distributions, or training introduces further communication
and numerical constraints. More efficient inference designs may reduce energy
per task; deployment scale and manufacturing costs determine their broader
resource impact.
